\section{Scope: what does not separate}
\label{sec:scope}

Two further results bound the claim, and we report both at the same
evidentiary standard as the wins.

\begin{table}[t]
\centering
\small
\tabletasktwo
\caption{Multi-hop variant, nine seeds per compared arm, accuracy at the
matched final checkpoint (chance 0.031, demonstration bar 0.094; readings
above the bar in bold). Three of nine matrix fast-weight seeds learn the
task; zero of nine flat-vector seeds do. The pooled paired interval
straddles zero and is additionally flagged non-decision-grade by a
pre-registered batch-effect gate, so this table is a trainability
disclosure, not a capability claim.}
\label{tab:task2}
\end{table}

\textbf{Multi-hop recall separates only as seed variance.} On the
multi-hop variant (compositional queries over a single Hamiltonian
$K$-cycle, trained at hop depths 1--2), the contender clears the
demonstration bar in 3 of 9 seeds (0.334, 0.479, 0.391) and reads chance
in the other six; the flat-vector arm reads chance in all
nine. % <!-- evidence: C7 -->
The registered adjudication rule scores this as
trainability/seed-variance: the hypothesis that multi-hop recall is a hard
capability boundary for this architecture at this scale is rejected (some
seeds reach it), but the pooled paired interval $(-0.020,\, 0.268)$
straddles zero and the pre-registered batch-effect gate flags the pooled
cohort (variance ratio 6.14 against a 4.0 bound), so no capability
separation is claimed in either direction. % <!-- evidence: C7 -->
The bar-clearing seeds also do not generalize: held-out hop depths 3--4
read chance, and the first bar-clearing seed's partial recall collapses to
0.010 when its episode is embedded at any extended
horizon. % <!-- evidence: C7 -->
Multi-hop composition at this scale is an open training problem, not a
demonstrated capability.

\begin{table}[t]
\centering
\small
\tablekstress
\caption{Stress cell at $K{=}48$ bindings on the same $d{=}64$ state
($K/d = 0.75$; chance 0.021). Every arm, including the matrix fast-weight
model, reads chance: the recall separation of Table~\ref{tab:main} is
load-bounded and does not extend to this operating point at this scale
and budget.}
\label{tab:k48}
\end{table}

\textbf{The separation is load-bounded.} At $K = 48$ bindings on the same
64-dimensional state ($K/d = 0.75$), all three arms read chance: 0.0189
(contender), 0.0195 (flat-vector), 0.0218
(transformer). % <!-- evidence: C8 -->
This cell was registered as locate-only, so it carries no tier
arithmetic; what it establishes is that the ceiling-level recall at
$K/d = 0.5$ does not come for free at higher load under the same training
budget. Where between the two load points the capability degrades, and
whether more training moves the boundary, are open questions this paper
does not answer.
